\documentclass{stylefiles/capstone}
\usepackage{listings}
\usepackage{url}
\usepackage{hyperref}
\usepackage{booktabs}
\usepackage{multirow}
\usepackage[numbers]{natbib}

\bibliographystyle{plainnat}

\begin{document}

\title{AI-Powered Smart Contract Bytecode Decompiler}

\capstoneyear{2025}
\capstonedocument{Seminar}
\capstonesemester{Spring}

\author{Karim Ali}
\affiliation{%
  \institution{Computer Science, NYUAD}
}
\email{karim.ali@nyu.edu}

\author{Eyerusalem Hawoltu Afework}
\affiliation{%
  \institution{Computer Science, NYUAD}
}
\email{eh3115@nyu.edu}

\author{Tamer Abdelaziz}
\affiliation{%
  \institution{Computer Science, NYUAD}
}
\email{tamer.abdelaziz@nyu.edu}

\renewcommand{\shortauthors}{Ali, Afework, Abdelaziz}

\begin{abstract}
Less than one percent of deployed Ethereum smart contracts have publicly verified
source code, leaving most on-chain logic opaque and difficult to audit. This
project builds an end-to-end pipeline that recovers high-level Solidity from
low-level Ethereum Virtual Machine (EVM) bytecode by way of an intermediate
control-flow-graph (CFG) representation: \emph{bytecode $\rightarrow$ CFG
$\rightarrow$ Solidity}, where the final stage is performed by a fine-tuned large
language model (LLM).

To support this study, we assemble a large-scale dataset by cloning the
Smart-Contract-Sanctuary-Ethereum repository (tintinweb), which contains hundreds
of thousands of real-world mainnet contracts. From this corpus we process
$292{,}837$ contract files, extract their deployed bytecode, and reconstruct
control-flow graphs using a substantially extended version of
\texttt{evm\_cfg\_builder}, pairing each CFG with its corresponding Solidity
function definition.

A central contribution of this work is the engineering and evaluation of the
\emph{bytecode-to-CFG} front-end. The upstream \texttt{evm\_cfg\_builder} is
unmaintained and fails to recover function boundaries on modern compiler output.
We add a breadth-first dispatcher-reconstruction pass, a merged
$453{,}208$-entry selector$\rightarrow$signature table, robustness fixes for
Solidity~0.8.x bytecode, and support for the Shanghai \texttt{PUSH0} opcode. On a
held-out evaluation of $109{,}679$ contracts spanning $21$ Solidity versions
(0.8.0--0.8.20), these changes raise function-name-recovery $F_1$ from $0.338$ to
$0.673$ (precision $0.888$, recall $0.648$) on the contracts whose bytecode can
contain recoverable functions, while precision \emph{rises} alongside recall,
confirming that the gains reflect correct recovery rather than metric inflation.

For the CFG-to-Solidity stage we adapt \textsc{Nova}, an assembly-to-C
generative model with hierarchical attention, to the Solidity domain and
fine-tune it on our aligned dataset. Preliminary evaluation on held-out functions
yields a BLEU-4 of $0.833$ and an exact-match rate of $70.8\%$. This work
provides one of the first systematic, mainnet-scale studies of CFG-driven,
learned smart-contract decompilation, and lays the groundwork for hybrid pipelines
that combine static analysis with LLM reasoning.
\end{abstract}

\keywords{Ethereum smart contracts, EVM bytecode, decompilation, control-flow
graphs, large language models, Nova, hierarchical attention, selector resolution,
PUSH0, blockchain security, Solidity reconstruction, bytecode analysis}

\maketitle

% ===========================================================================
\section{Introduction}

Decentralized applications (dApps) and decentralized finance (DeFi) rely heavily
on smart contracts, which collectively secure and move billions of dollars across
public blockchains. Despite this scale, more than 99\% of deployed Ethereum
contracts lack publicly verified source code~\cite{etherscan}. Once compiled into
Ethereum Virtual Machine (EVM) bytecode, a contract loses its high-level
structure, type information, and developer intent, leaving auditors, developers,
and users with only low-level opcodes to reason about. This opacity severely
limits transparency and makes security analysis, incident response, and forensic
investigations significantly more difficult.

Traditional attempts to reverse-engineer smart contracts rely on
static-analysis–based decompilation tools such as Dedaub~\cite{neville2023dedaub}
and Panoramix~\cite{panoramix}. Although these systems recover partial
control-flow structure, they often produce incomplete, noisy, or misleading
output—particularly for modern optimized bytecode or contracts involving complex
financial logic. Prior work suggests that static approaches alone are
fundamentally limited in their ability to reconstruct high-level semantics from
EVM execution patterns~\cite{darwish2023bytecode}. These limitations motivate
exploring learning-based approaches that can generalize beyond hand-crafted
heuristics.

Recent advances in large language models (LLMs) have demonstrated strong
capabilities in code understanding, translation, and synthesis, exemplified by
GPT-3~\cite{brown2020language} and CodeT5+~\cite{wang2023codet5plus}. Most
relevant to this work is \textsc{Nova}~\cite{jiang2025nova}, a model that
decompiles low-level x86 assembly into C using a \emph{hierarchical attention}
mechanism designed for the low information density and long-range dependencies of
machine code. The parallel between Nova's assembly-to-C task and our
bytecode-to-Solidity task is exact, motivating us to adapt Nova to the Solidity
domain rather than evaluate general-purpose chat models.

This project therefore frames decompilation as a two-stage problem. First, the
deployed runtime bytecode is lifted into a structured, textual control-flow graph
(CFG) that exposes basic blocks, dispatcher structure, and per-function selectors.
Second, a fine-tuned Nova model translates the CFG into a Solidity function. To
enable both stages we construct a large-scale dataset by cloning the
Smart-Contract-Sanctuary-Ethereum repository~\cite{smart_contract_sanctuary},
which archives hundreds of thousands of real mainnet contracts. From this corpus
we curate a filtered subset of $292{,}837$ contract files by removing legacy
compiler versions and eliminating contracts with unresolved dependencies. For
each remaining contract we extract all available Solidity function definitions and
pair each function with a CFG reconstructed from its compiled bytecode using an
extended version of \texttt{evm\_cfg\_builder}~\cite{evm_cfg_builder}.

Compilation is inherently lossy: variable names, comments, type annotations, and
higher-level control structures are removed during the source-to-bytecode
transformation. Recovering these abstractions requires both an accurate CFG
front-end and a model capable of understanding opcode patterns, long-range
dependencies, and control-flow semantics. Our contributions are:

\begin{enumerate}
  \item a scalable pipeline that aligns real mainnet Solidity functions with CFGs
        reconstructed from their bytecode;
  \item a substantially extended \texttt{evm\_cfg\_builder} with a breadth-first
        dispatcher-reconstruction pass, a merged $453{,}208$-entry selector
        table, Solidity~0.8.x robustness fixes, and \texttt{PUSH0} opcode support,
        together with a controlled, version-stratified evaluation on $109{,}679$
        contracts; and
  \item a Solidity adaptation of the Nova architecture, fine-tuned on the aligned
        dataset, with preliminary end-to-end decompilation results.
\end{enumerate}

% ===========================================================================
\section{Related Work}

Several research efforts have explored smart contract analysis, decompilation, and
the use of machine learning for code understanding. This section highlights
related work across three domains: smart contract decompilers, learning-based code
models, and the Nova architecture our pipeline builds on.

Dedaub~\cite{neville2023dedaub} is a leading smart contract decompiler that
reconstructs high-level contract logic from EVM bytecode using static analysis and
control-flow-graph recovery. It produces human-readable pseudo-code and has been
widely used in both research and industry. However, its rule-based design limits
its adaptability to previously unseen or obfuscated bytecode structures.
Similarly, Panoramix~\cite{panoramix} disassembles contracts into a
pseudo-assembly form and reconstructs function boundaries using symbolic execution
and jump resolution. While effective in certain cases, its output lacks
Solidity-like readability and fails on many edge cases or newer compiler patterns.

More recently, researchers have explored using LLMs to analyze smart contracts.
Darwish~\cite{darwish2023bytecode} evaluated GPT-3.5, GPT-4, and CodeLlama for
detecting vulnerabilities in contracts decompiled with
Dedaub~\cite{neville2023dedaub}, demonstrating the potential of language models for
security analysis but not attempting to reconstruct source from raw bytecode. Our
project differs in that we train a model directly on aligned bytecode-derived CFGs
and source code, eliminating the dependency on external decompilers.

The foundational work on LLMs comes from Brown et al.~\cite{brown2020language}, who
introduced GPT-3 and demonstrated strong few-shot learning and code generation.
This inspired a wave of code-specialized models, of which
CodeT5+~\cite{wang2023codet5plus} is a notable open-source example trained for
program understanding, synthesis, and translation. The model most directly
relevant to our task is \textsc{Nova}~\cite{jiang2025nova}, which targets
\emph{binary} decompilation (assembly~$\rightarrow$~C). Nova introduces a
hierarchical attention mechanism—where some attention heads attend to individual
instructions while others attend to a learned summary of instruction blocks—and a
contrastive pre-training objective, allowing it to outperform substantially larger
general-purpose models on assembly decompilation. We adopt Nova's architecture
unchanged and re-target it from assembly/C to CFG/Solidity.

Finally, Etherscan~\cite{etherscan} plays a central role in the Ethereum ecosystem
as a public block explorer and verification service. Our dataset builds on the
Smart-Contract-Sanctuary-Ethereum repository~\cite{smart_contract_sanctuary}, which
archives mainnet contracts and their verified sources at scale, allowing us to
construct aligned bytecode--source code pairs.

% ===========================================================================
\section{Methodology}

This section describes the full pipeline used to construct the dataset, extract
control-flow graphs (CFGs) from EVM bytecode, and fine-tune the Nova model to
reconstruct Solidity. The methodology consists of four stages: dataset
construction, CFG extraction (and the engineering required to make it reliable),
function-level alignment, and model fine-tuning.

\subsection{Overview}

The overall workflow proceeds as follows:

\begin{enumerate}
    \item Collect verified Solidity contracts from a large public
          repository~\cite{smart_contract_sanctuary}.
    \item Group contracts by compiler version and compile them to bytecode.
    \item Extract function-level CFGs using an extended version of
          \texttt{evm\_cfg\_builder}~\cite{evm_cfg_builder}.
    \item Align each CFG with its corresponding Solidity source function.
    \item Fine-tune the Nova model to translate a CFG into a Solidity function.
    \item Evaluate the pipeline at two levels: function-name recovery in the CFG
          front-end, and end-to-end source reconstruction with
          BLEU~\cite{papineni2002bleu} and CodeBERT~\cite{feng2020codebert}.
\end{enumerate}

\subsection{Dataset Construction}

\subsubsection{Collecting Contracts and Grouping by Version}

Dataset construction begins with a large-scale clone of the Smart Contract
Sanctuary Ethereum repository~\cite{smart_contract_sanctuary}, which hosts hundreds
of thousands of real-world mainnet contracts.

\begin{verbatim}
contracts/mainnet/
 |-- 000/
 |-- 001/
 |-- ...
\end{verbatim}

All Solidity source files (\texttt{*.sol}) are extracted, cleaned, and grouped by
compiler version based on their \texttt{pragma solidity} declarations. Each
compiler version is placed into its own directory:

\begin{verbatim}
Contracts_By_Version/
 |-- 0.8.0/
 |-- 0.8.4/
 |-- 0.8.7/
 |-- 0.8.9/
 |-- ...
\end{verbatim}

After filtering out unsupported legacy versions and removing contracts with
unresolved dependencies, the final dataset contains \textbf{292{,}837} Solidity
contract files spanning multiple minor versions of the compiler.

\subsubsection{Compiling Contracts}

To reproduce the deployed bytecode as faithfully as possible, each contract is
compiled using the exact Solidity compiler version extracted from its pragma. This
is performed using \texttt{solc-select}, falling back to the closest compatible
patch version when needed. Compilation yields runtime bytecode files
(\texttt{*.hex}), compilation logs, and failure reports, and is driven by a
resumable pipeline that logs malformed files and timeouts so that the full corpus
can be processed without manual restarts.

\subsection{CFG Extraction and \texttt{evm\_cfg\_builder} Enhancements}
\label{sec:cfg-extract}

CFGs are extracted using \texttt{evm\_cfg\_builder}~\cite{evm_cfg_builder}, a
static analysis tool originally provided by Trail of Bits. For each bytecode file
the tool outputs a dispatcher graph (mapping selectors to code blocks), a full CFG,
and a serialized textual form listing basic blocks and opcodes. Each CFG captures
fundamental control-flow information such as \texttt{JUMP}, \texttt{JUMPI},
function entry points, revert blocks, and calldata loaders; these structural
elements form the input to the downstream model.

\paragraph{Why the upstream tool is insufficient.}
The upstream project has not been actively maintained in recent years and was
written for older compiler output. On modern Solidity~0.8.x bytecode it (i)
crashes when a jump target lands inside \texttt{PUSH} data bytes, (ii) fails to
compile its built-in $417{,}765$-entry selector table under Python~3.10, (iii)
relies solely on value-set analysis (VSA), which leaves many reachable basic
blocks unanalyzed and lets control flow ``bleed'' across function boundaries, and
(iv) cannot decode the Shanghai \texttt{PUSH0} opcode emitted by Solidity~$\geq$
0.8.20. These failures motivated a series of targeted changes to the tool, which
we describe below and evaluate in Section~\ref{sec:cfg-eval}.

\paragraph{Package re-export.}
The cloned repository moved the \texttt{CFG} class into
\texttt{cfg/cfg.py} but left \texttt{cfg/\_\_init\_\_.py} empty. We re-export
\texttt{CFG}, \texttt{BasicBlock}, and \texttt{Function} so that
\texttt{from evm\_cfg\_builder.cfg import CFG} resolves correctly throughout the
pipeline.

\paragraph{Breadth-first dispatcher reconstruction.}
The core enhancement is a breadth-first traversal that recovers basic blocks the
VSA misses. It comprises four additions to \texttt{cfg/cfg.py}:
\begin{itemize}
  \item \texttt{\_static\_jump\_target(bb)} scans backwards through a block's
        instructions to find the last \texttt{PUSH} and returns the corresponding
        target \texttt{BasicBlock} (or \texttt{None}).
  \item \texttt{\_scan\_selector\_dispatch()} scans all basic blocks for the
        canonical pattern
        \texttt{PUSH4 <selector>; EQ; PUSH2 <body\_pc>; JUMPI} and returns a map
        \{selector~$\rightarrow$~body\_pc\}. This works for both inline dispatchers
        and the shared ABI-decoder layout emitted by Solidity~$\geq$ 0.8.x.
  \item \texttt{\_bfs\_blocks(entry, stop\_at, key)} performs a BFS from an entry
        block, preferring function-specific VSA edges first (to avoid
        cross-function bleeding) and falling back to static \texttt{PUSH}
        resolution only for blocks the VSA never reached, descending the
        \texttt{JUMPI} false branch when resolving statically.
  \item \texttt{enhance\_cfgs\_with\_bfs()} is called after \texttt{CFG}
        construction; for each function it BFS-traverses from both the entry stub
        and the actual body (located via the dispatch scan), merges the results,
        and upgrades the function only when strictly more blocks are recovered than
        the VSA produced.
\end{itemize}

\paragraph{Selector-to-name resolution.}
Function dispatch is keyed by a four-byte selector
$\mathrm{sel}(f)=\texttt{keccak256}(\text{signature}(f))[0{:}4]$, so recovering
function \emph{names} requires a selector$\rightarrow$signature table. We began
with a curated list of $55{,}674$ known pairs ($55{,}668$ unique selectors)
derived from on-chain frequency data, e.g.\
\verb|0x70a08231 -> balanceOf| and \verb|0xa9059cbb -> transfer|. We discovered
that the disassembler loads its table ``JSON-first, no-merge'': because the
curated table exists, the tool's own $417{,}765$-entry built-in table was ignored
entirely, so the curated table had silently \emph{replaced} rather than augmented
it. The two tables also differ in format—the curated table stores \emph{bare}
names (\texttt{transfer}) while the built-in table stores \emph{typed} signatures
(\texttt{transfer(address,uint256)}). We therefore normalize the built-in entries
to bare names and merge the two with curated priority, producing a combined
\textbf{$453{,}208$-entry} table ($+397{,}540$ over the curated set, with only $19$
genuine name overrides). An optional-import guard additionally catches the
\texttt{OverflowError} raised when Python~3.10 cannot compile the large built-in
table, falling back gracefully so that CFG extraction still succeeds.

\paragraph{\texttt{PUSH0} opcode support.}
The bundled disassembler's default fork is \texttt{istanbul} (2019), whose
instruction table predates the Shanghai hard fork and therefore decodes
\verb|0x5f| (\texttt{PUSH0}) as \texttt{INVALID}. Because \texttt{INVALID}
terminates a basic block, any \texttt{PUSH0} in the dispatcher corrupts function
recovery, leaving only the fallback. \texttt{PUSH0} became the default in
Solidity~0.8.20. We register \texttt{PUSH0} (zero operand bytes, pushes the
constant $0$) in every fork table at import, restoring correct disassembly.

\paragraph{Robustness fix.}
Finally, we guard all \texttt{self.\_basic\_blocks[x]} lookups in
\texttt{compute\_functions} that could raise a \texttt{KeyError} on
Solidity~0.8.x when a jump target lands in the middle of \texttt{PUSH} data bytes,
eliminating a class of crashes that previously aborted extraction.

\paragraph{Pipeline integration.}
In \texttt{nova/bytecode\_to\_cfg.py}, the local extended \texttt{evm\_cfg\_builder}
clone is prepended to \texttt{sys.path} at startup (shadowing any pip-installed
copy on both macOS and the HPC cluster), \texttt{enhance\_cfgs\_with\_bfs()} is
invoked immediately after \texttt{CFG(bytecode)} construction, and the enhanced
\texttt{func.basic\_blocks} are used directly, with the older manual BFS retained
only as an emergency fallback.

\subsection{Function-Level Alignment}

For each contract, all function definitions are parsed from the Solidity source and
matched to their CFG representation based on function selectors and dispatch
structure. Each aligned entry takes the following JSON form:

\noindent\begin{small}
\begin{verbatim}
{
  "withdrawForeignERC721(address,uint256)": {
    "solidity_definition":
         "function withdrawForeignERC721(...) {...}",
    "cfg_representation":
         "Function withdrawForeignERC721.\nBasic Blocks: ..."
  }
}
\end{verbatim}
\end{small}

A normalized JSONL dataset is also created, where each line is a single function
instance:

\noindent\begin{small}
\begin{verbatim}
{
  "version": "0.8.9",
  "function_name": "toggleSwap(bool)_v175",
  "cfg_representation": "...",
  "solidity_definition":
        "function toggleSwap(bool _swapEnabled) ..."
}
\end{verbatim}
\end{small}

The suffix (e.g., \texttt{\_v175}) uniquely distinguishes functions with identical
signatures across different contracts.

\subsection{Decompilation Task Definition}

The core task is:

\begin{quote}
\textit{Given only a textual CFG generated from EVM bytecode, reconstruct the most
likely Solidity implementation of the function.}
\end{quote}

A reconstruction must preserve the function signature, infer visibility, modifiers,
and revert conditions, reflect the control flow implied by the CFG, and resemble
the ground-truth Solidity implementation. This framing mirrors a realistic
reverse-engineering scenario in which only runtime bytecode is available.

\subsection{Nova Model and Fine-Tuning}
\label{sec:nova}

For the CFG-to-Solidity stage we adapt \textsc{Nova}~\cite{jiang2025nova}. We keep
Nova's architecture entirely unchanged—\texttt{modeling\_nova.py} (the
\texttt{NovaForCausalLM} model, \texttt{NovaAttention} hierarchical attention, and
\texttt{NovaTokenizer}) and \texttt{generation\_utils.py} (the generation logic)
are used as released—because the parallel between Nova's assembly-to-C task and our
CFG-to-Solidity task is exact. Our work consists of re-targeting the data,
training, and evaluation around this architecture:

\begin{itemize}
  \item \textbf{Dataset preparation} (\texttt{prepare\_solidity\_dataset.py}, new).
        We read the aligned CFG/Solidity function pairs and insert
        \texttt{<label-N>} tokens after each EVM opcode line (\texttt{JUMPDEST},
        \texttt{PUSH2}, \texttt{CALLDATALOAD}, \dots), which is what enables Nova's
        hierarchical attention to form per-block summaries. The Solidity compiler
        version is included in the prompt so the model conditions on it, and all
        versions are combined into a single dataset where the version string acts
        as a conditioning signal. The prompt template is
        \texttt{"\# This is the EVM CFG for a Solidity \{version\} function:\textbackslash n\{cfg\}\textbackslash nWhat is the Solidity source code?"}.
  \item \textbf{On-the-fly attention mask} (\texttt{dataset.py}, modified to add
        \texttt{SolidityDataset}). Nova pre-computes a per-sample
        \texttt{nova\_attention\_mask} of size $2048\times2048$. Materializing this
        for the full training set would require on the order of a terabyte of
        storage; we instead build the mask per sample during training, eliminating
        the storage bottleneck.
  \item \textbf{Fine-tuning} (\texttt{finetune\_solidity.py}, new). We load
        \texttt{deepseek-ai/deepseek-coder-1.3b-base} into
        \texttt{NovaForCausalLM} and fine-tune with DeepSpeed ZeRO-2 in BF16 for
        $3$ epochs.
  \item \textbf{Inference and evaluation} (\texttt{inference\_solidity.py} and
        \texttt{evaluate\_solidity.py}, new). Inference takes an EVM CFG and emits
        predicted Solidity (single-function and batch modes); evaluation computes
        BLEU-1 through BLEU-4, CodeBERT cosine similarity, and exact-match rate.
  \item \textbf{Job orchestration} (\texttt{hpc/}). SLURM scripts drive dataset
        preparation, multi-GPU training, and inference/evaluation on the cluster.
\end{itemize}

% ===========================================================================
\section{CFG Front-End Evaluation}
\label{sec:cfg-eval}

Because the quality of the entire pipeline is bounded by the quality of the CFG
front-end, we evaluate the front-end independently before training the model. We
measure how accurately the enhanced \texttt{evm\_cfg\_builder} recovers the
\emph{public function names} of a contract from its bytecode alone, using the
verified Solidity source as ground truth.

\subsection{Setup and Metrics}

We evaluate on \textbf{$109{,}679$ contracts} across \textbf{$21$ Solidity
versions} (0.8.0--0.8.20), each providing deployed bytecode and cleaned source.
Evaluation runs on the HPC cluster via SLURM with $16$ parallel worker processes,
processing the full set in under ten minutes and emitting one auditable row per
contract.

Let $\widehat{F}$ be the set of named functions recovered from the bytecode and
$F$ the set of public/external functions parsed from the source, with
$M=\widehat{F}\cap F$. We report precision $|M|/|\widehat{F}|$, recall $|M|/|F|$,
their harmonic mean $F_1$, exact set match $\mathbb{1}[\widehat{F}=F]$, and
sentence-level BLEU over the sorted name lists. Precision answers ``when the tool
names a function, is it correct?''; recall answers ``what fraction of deployed
public functions are recovered?''.

\subsection{Ablation}

We isolate each change as a separate configuration. Two of the changes are
\emph{measurement-correctness} fixes that stop scoring the extractor on functions
that cannot be present in the deployed bytecode by construction, and two are
genuine \emph{capability} fixes.

\paragraph{E1: Interface exclusion (measurement).}
A flattened source file contains the deployed contract plus every imported
interface (e.g.\ \texttt{IUniswapV2Router02}). Interface functions are never
compiled into the contract's own bytecode—they are external ABIs used to call
\emph{other} contracts—so counting them as deployed public functions inflated the
recall denominator. Excluding \texttt{interface} blocks (and stripping comments to
avoid spurious matches) raised recall from $0.290$ to $0.383$ with precision
essentially unchanged.

\paragraph{E2: Merged selector table (capability).}
Merging the curated and built-in selector tables into the combined
$453{,}208$-entry table (Section~\ref{sec:cfg-extract}) resolved additional
selectors, trading a small precision cost for a recall gain (net $F_1$ positive).

\paragraph{E3: Library-stub exclusion (measurement).}
A large fraction of ``contracts'' are in fact deployed Solidity \emph{libraries}
whose runtime is an $86$-byte delegatecall guard
(\texttt{PUSH20 <self-address>; ADDRESS; EQ; \dots}) containing no callable
functions by construction—library functions are \texttt{internal} and inlined into
callers. Such stubs constitute $58.8\%$ of 0.8.0 and $83.1\%$ of 0.8.1, and
scoring them against a full multi-function source file is meaningless. A bytecode
classifier labels each contract \texttt{runtime}, \texttt{library\_stub},
\texttt{tiny}, or \texttt{empty}; excluding the $35{,}530$ non-runtime contracts
yields the fair, runtime-only headline.

\paragraph{E4: \texttt{PUSH0} opcode support (capability).}
Registering \texttt{PUSH0} (Section~\ref{sec:cfg-extract}) restores function
recovery on Shanghai-era bytecode. A key empirical finding is that the version
buckets are defined by source \emph{pragma}, not by the compiler actually used at
deployment: $88.2\%$ of ``0.8.0-pragma'' runtime contracts contain \texttt{PUSH0},
so the single disassembler fix improves \emph{every} version. A per-contract audit
confirms the gains are correct recovery: contracts that previously produced
\emph{empty} CFGs now recover full, correctly-named function sets at precision
$\approx 1.0$.

\begin{table}[t]
\centering
\caption{All-contracts weighted ablation ($109{,}679$ contracts).}
\label{tab:ablation-all}
\small
\begin{tabular}{l c c c}
\toprule
Configuration & Precision & Recall & $F_1$ \\
\midrule
V1: Baseline                    & 0.5647 & 0.2895 & 0.3378 \\
V2: + Interface exclusion       & 0.5635 & 0.3834 & 0.4106 \\
V3: + Merged $453$k table       & 0.5554 & 0.4024 & 0.4168 \\
V4: + \texttt{PUSH0} fix        & \textbf{0.6003} & \textbf{0.4378} & \textbf{0.4546} \\
\bottomrule
\end{tabular}
\end{table}

\begin{table}[t]
\centering
\caption{Runtime-only weighted headline (library stubs / empty / tiny bytecode
excluded; $73{,}709$ contracts).}
\label{tab:ablation-rt}
\small
\begin{tabular}{l c c c c}
\toprule
Configuration & Precision & Recall & $F_1$ & Exact \\
\midrule
V3$_{\text{rt}}$: merged, no stubs & 0.8217 & 0.5953 & 0.6167 & 0.097 \\
V4$_{\text{rt}}$: + \texttt{PUSH0} & \textbf{0.8881} & \textbf{0.6479} & \textbf{0.6726} & \textbf{0.109} \\
\bottomrule
\end{tabular}
\end{table}

\subsection{Results}

Table~\ref{tab:ablation-all} reports the weighted headline over all contracts, and
Table~\ref{tab:ablation-rt} the fair, runtime-only headline. On the $73{,}709$
contracts whose bytecode can contain recoverable functions, the final
configuration attains precision $0.888$, recall $0.648$, and $F_1=0.673$, up from
a baseline $F_1$ of $0.338$. Crucially, precision \emph{rises} from $0.56$ to
$0.89$ as recall rises: had any stage simply emitted more guesses, precision would
have fallen, so the increase indicates the additional recovered names are correct.
Table~\ref{tab:perversion} gives the final per-version breakdown; precision is
uniformly high ($0.68$--$0.93$), with recall the limiting axis, reflecting
inherited base-class functions whose dispatch is not always followed.

\begin{table}[t]
\centering
\caption{Final runtime-only results per version. $n$ is runtime contracts; $s$ is
excluded library stubs.}
\label{tab:perversion}
\scriptsize
\begin{tabular}{l r r c c c c}
\toprule
Version & $n$ & $s$ & Prec. & Rec. & $F_1$ & BLEU \\
\midrule
0.8.0  & 14{,}749 & 21{,}305 & 0.895 & 0.775 & 0.802 & 0.530 \\
0.8.1  & 283      & 1{,}416  & 0.800 & 0.717 & 0.718 & 0.419 \\
0.8.2  & 431      & 482      & 0.682 & 0.815 & 0.697 & 0.326 \\
0.8.3  & 404      & 340      & 0.795 & 0.730 & 0.711 & 0.427 \\
0.8.4  & 9{,}791  & 2{,}499  & 0.926 & 0.657 & 0.675 & 0.475 \\
0.8.5  & 1{,}006  & 239      & 0.830 & 0.896 & 0.831 & 0.565 \\
0.8.6  & 968      & 362      & 0.734 & 0.606 & 0.587 & 0.277 \\
0.8.7  & 5{,}884  & 1{,}005  & 0.905 & 0.611 & 0.643 & 0.413 \\
0.8.8  & 270      & 93       & 0.891 & 0.719 & 0.728 & 0.464 \\
0.8.9  & 14{,}420 & 1{,}962  & 0.894 & 0.600 & 0.621 & 0.365 \\
0.8.10 & 6{,}491  & 895      & 0.885 & 0.550 & 0.608 & 0.214 \\
0.8.11 & 1{,}323  & 370      & 0.843 & 0.604 & 0.637 & 0.257 \\
0.8.12 & 968      & 701      & 0.852 & 0.636 & 0.647 & 0.356 \\
0.8.13 & 2{,}065  & 745      & 0.822 & 0.539 & 0.591 & 0.239 \\
0.8.14 & 811      & 242      & 0.880 & 0.583 & 0.616 & 0.307 \\
0.8.15 & 2{,}140  & 642      & 0.867 & 0.610 & 0.642 & 0.319 \\
0.8.16 & 2{,}961  & 477      & 0.892 & 0.648 & 0.687 & 0.402 \\
0.8.17 & 6{,}504  & 1{,}077  & 0.892 & 0.593 & 0.614 & 0.370 \\
0.8.18 & 540      & 96       & 0.798 & 0.709 & 0.672 & 0.301 \\
0.8.19 & 1{,}133  & 193      & 0.894 & 0.625 & 0.654 & 0.390 \\
0.8.20 & 567      & 62       & 0.902 & 0.634 & 0.648 & 0.423 \\
\midrule
\textbf{Wtd.} & \textbf{73{,}709} & \textbf{35{,}530} & \textbf{0.888} & \textbf{0.648} & \textbf{0.673} & \textbf{0.400} \\
\bottomrule
\end{tabular}
\end{table}

% ===========================================================================
\section{End-to-End Decompilation with Nova}
\label{sec:nova-eval}

With a reliable CFG front-end in place, we fine-tune the Nova model
(Section~\ref{sec:nova}) on the aligned dataset and evaluate end-to-end
CFG-to-Solidity reconstruction. The model is trained on $315{,}794$ training and
$5{,}000$ validation function pairs; reconstruction quality is measured with
BLEU~\cite{papineni2002bleu} for syntactic similarity, CodeBERT~\cite{feng2020codebert}
cosine similarity for semantic alignment, and an exact-match rate.

\subsection{Why BLEU and CodeBERT}
BLEU measures surface-form n-gram overlap and is sensitive to variable renaming,
formatting, and benign reordering—so it under-credits logically correct but
syntactically different completions. CodeBERT cosine similarity,
$\text{sim}(g,\hat{g}) = E(g)\!\cdot\! E(\hat{g}) \,/\, (\lVert E(g)\rVert\,\lVert E(\hat{g})\rVert)$,
where $E(\cdot)$ is the CodeBERT encoder, captures semantic alignment that BLEU
misses. Exact string matching or AST equivalence is too strict because compilation
is lossy, and full symbolic equivalence requires executable contract context that
is unavailable for most entries; BLEU plus CodeBERT therefore provides a scalable,
reproducible, and meaningful measure.

\subsection{Preliminary Results}
Table~\ref{tab:nova} reports a preliminary evaluation of the fine-tuned model on
$32{,}079$ held-out functions. The model reconstructs the original function
\emph{verbatim} in $70.8\%$ of cases and achieves a BLEU-4 of $0.833$; $22{,}817$
reconstructions score BLEU~$\geq 0.9$. Qualitatively, the model reliably reproduces
ownership and configuration functions and common ERC-20 patterns, while more
complex behaviours involving loops over dynamic arrays, inter-contract calls, or
multi-step invariants remain harder. These results are preliminary: a full retrain
and an evaluation over the complete held-out set are in progress and will be
incorporated in the final version of this report.

\begin{table}[t]
\centering
\caption{Preliminary end-to-end reconstruction results for the fine-tuned Nova
model ($32{,}079$ held-out functions).}
\label{tab:nova}
\small
\begin{tabular}{l c}
\toprule
Metric & Value \\
\midrule
Functions evaluated         & $32{,}079$ \\
Avg.\ BLEU-4                 & $0.833$ \\
Exact match                  & $70.8\%$ \;($22{,}716/32{,}079$) \\
BLEU $\geq 0.9$              & $22{,}817$ \\
BLEU $< 0.5$                 & $4{,}747$ \\
\bottomrule
\end{tabular}
\end{table}

% ===========================================================================
\section{Project Timeline}

\begin{enumerate}

\item \textbf{Literature Review and Problem Definition (Week 1--2):}
Surveyed prior work on EVM bytecode analysis, decompilation, and smart contract
security—including Dedaub~\cite{neville2023dedaub} and Panoramix~\cite{panoramix}—
and learning-based code models including GPT-3~\cite{brown2020language},
CodeT5+~\cite{wang2023codet5plus}, and the assembly-to-C model
Nova~\cite{jiang2025nova}. This clarified the research gap and established the
central question of CFG-driven, learned decompilation.

\item \textbf{Dataset Acquisition (Week 3):}
Cloned the Smart Contract Sanctuary repository and collected $292{,}837$ Solidity
files across many compiler versions, extracting pragma metadata for version-based
grouping.

\item \textbf{Bytecode Compilation and CFG Extraction (Week 4--6):}
Compiled all contracts with \texttt{solc-select} and fallback version resolution,
producing runtime bytecode, and extracted CFGs using the extended
\texttt{evm\_cfg\_builder}. Failures and timeouts were logged with a resumable
pipeline.

\item \textbf{CFG Front-End Engineering and Evaluation (Week 6--8):}
Implemented and evaluated the BFS dispatcher reconstruction, merged
$453{,}208$-entry selector table, Solidity~0.8.x robustness fixes, and
\texttt{PUSH0} support, running the version-stratified evaluation on $109{,}679$
contracts (Section~\ref{sec:cfg-eval}).

\item \textbf{Function-Level Alignment and Dataset Assembly (Week 8--9):}
Aligned each CFG with its Solidity function and produced the JSON/JSONL datasets,
including the Nova-specific \texttt{<label-N>} tokenization.

\item \textbf{Nova Adaptation and Fine-Tuning Setup (Week 9--11):}
Adapted Nova to Solidity (new dataset-prep, on-the-fly attention masks, training,
inference, and evaluation scripts) and configured multi-GPU DeepSpeed training on
the cluster.

\item \textbf{Training and Evaluation (Week 11--13):}
Fine-tuned the model and evaluated reconstructions with
BLEU~\cite{papineni2002bleu} and CodeBERT~\cite{feng2020codebert}; a full retrain
and complete-set evaluation are in progress.

\item \textbf{Final Report and Presentation (Week 13--14):}
Refined the report, verified citations, prepared slides, and produced the final
PDF.

\end{enumerate}

% ===========================================================================
\section{Resources and Budget}

The project relied primarily on NYUAD's Jubail high-performance computing cluster,
which provided the CPU resources for large-scale CFG extraction (sixteen-worker
parallel jobs over $109{,}679$ contracts) and the GPU resources for Nova
fine-tuning and inference, at no direct monetary cost to the student. Because the
final pipeline performs decompilation with a self-hosted, fine-tuned Nova model, it
does not depend on any paid commercial LLM API at inference time. Software
dependencies (\texttt{solc-select}, \texttt{evm\_cfg\_builder}, PyTorch, DeepSpeed,
and the Nova codebase) are open source.

% ===========================================================================
\section{Conclusion}

This project builds an end-to-end pipeline that recovers high-level Solidity from
EVM bytecode through an intermediate control-flow-graph representation. To enable
it, we assembled a mainnet-scale dataset of $292{,}837$ real contracts, compiled
them, and reconstructed function-level CFGs. The central engineering contribution
is a substantially extended \texttt{evm\_cfg\_builder}: a breadth-first dispatcher
reconstruction pass, a merged $453{,}208$-entry selector$\rightarrow$signature
table, Solidity~0.8.x robustness fixes, and support for the Shanghai
\texttt{PUSH0} opcode. A controlled, version-stratified evaluation over $109{,}679$
contracts shows these changes raise function-name-recovery $F_1$ from $0.338$ to
$0.673$ on extractable contracts—with precision rising from $0.56$ to $0.89$—and
reveals that Solidity pragma is a poor proxy for the deployed compiler, so a single
\texttt{PUSH0} fix improves every version.

For the reconstruction stage we adapted the Nova architecture to Solidity, keeping
its hierarchical-attention model unchanged while rewriting the data, training, and
evaluation around it, and addressing a terabyte-scale attention-mask storage
bottleneck with on-the-fly mask construction. Preliminary results are encouraging:
the fine-tuned model reconstructs $70.8\%$ of held-out functions exactly with a
BLEU-4 of $0.833$. A full retrain and complete-set evaluation are underway and will
be reported in the final version.

Taken together, the results demonstrate that CFG-driven, learned decompilation is
both feasible and promising. The findings establish a foundational, reproducible
benchmark and point toward hybrid systems that combine static analysis with
LLM-based reasoning, advancing transparency, accessibility, and security in the
smart contract ecosystem.

\bibliography{refs}

\end{document}
\endinput
